\subsection{KissFeatures Message}

The \texttt{mrc/msg/KissFeatures} message encapsulates the extracted
geometric features for a single robot:

\begin{itemize}
    \item \texttt{keypoint\_xyz} (\texttt{float32[]}): 3D coordinates of the
          extracted keypoints stored as a flat, row-major array of the form
          $[x_0, y_0, z_0,\; x_1, y_1, z_1,\; \ldots]$, so the keypoint count
          is implicitly $|\texttt{keypoint\_xyz}| / 3$.
          This replaces the earlier \texttt{sensor\_msgs/PointCloud2} field,
          eliminating the PointCloud2 header and field-descriptor overhead
          ($\approx$13\,KB per message).

    \item \texttt{descriptors} (\texttt{uint16[]}): Flattened row-major array
          of quantised FasterPFH descriptors with
          $|\texttt{descriptors}| = (|\texttt{keypoint\_xyz}| / 3) \times
          \texttt{descriptor\_dim}$.
          Each value is quantised from its original \texttt{float32} range
          $[0, 100]$ to an unsigned 16-bit integer via
          \begin{equation}
              u = \mathrm{round}(v \times 655.35), \qquad
              v = u \;/\; 655.35,
          \end{equation}
          where the scale $655.35 = 65535/100$ maps the full \texttt{uint16}
          range $[0, 65535]$ onto $[0, 100]$.
          This reduces the descriptor array from 4\,bytes to 2\,bytes per
          element with a maximum quantisation error of
          $\varepsilon_{\max} = 0.5/655.35 \approx 7.6\times10^{-4}$ per bin.
          The FasterPFH histogram bins are normalised to sum to 100 per
          sub-histogram, so the value range fits entirely within
          \texttt{uint16} without clipping.

    \item \texttt{descriptor\_dim} (\texttt{uint32}): Dimensionality of each
          descriptor (33 for FasterPFH: three sub-histograms of 11 bins each).

    \item \texttt{descriptor\_type} (\texttt{string}): Identifier of the
          descriptor backend, e.g.\ \texttt{"FasterPFH"}.
\end{itemize}

\subsubsection*{Message Size}

Table~\ref{tab:kiss_size} reports measured wire sizes for a two-robot
proximity test (ec\_swift, front\,+\,back RGL lidar, $\approx$18\,700 raw
points per snapshot).

\begin{table}[h]
\centering
\caption{Measured \texttt{KissFeatures} wire size vs.\ raw lidar data
         ($\approx$3\,390 keypoints, descriptor dim\,=\,33)}
\label{tab:kiss_size}
\begin{tabular}{lrr}
\toprule
\textbf{Representation} & \textbf{Size} & \textbf{Ratio} \\
\midrule
Raw \texttt{PointCloud2} (2 lidars combined) & 292.9\,KB & $1.0\times$ (baseline) \\
KissFeatures --- original (\texttt{float32} descriptors, \texttt{PointCloud2} keypoints)
    & 490.0\,KB & $1.7\times$ \emph{larger} \\
\textbf{KissFeatures --- optimised (\texttt{uint16} descriptors, \texttt{float32[]} keypoints)}
    & \textbf{$\approx$258\,KB} & $\mathbf{1.1\times}$ \textbf{smaller} \\
\bottomrule
\end{tabular}
\end{table}

The dominant cost before optimisation was the descriptor array, which
accounted for 89\,\% of the total message size
($437$\,KB out of $490$\,KB).
Switching from \texttt{float32} to \texttt{uint16} encoding reduces the
descriptor array by 50\,\% ($437$\,KB\,$\to$\,$218$\,KB), yielding an overall
message reduction of 47\,\% and keeping KissFeatures below the raw
point-cloud bandwidth while providing a worst-case quantisation error of
$\varepsilon_{\max} \approx 7.6\times10^{-4}$ per bin.
